\Article{14}
Let us imagine the volume of the movable needle divided into infinitesimal elements. Let $x,y,z$ be the coordinates of an arbitrary element and $e$ the element of free magnetism it contains. On the vertical line through $G$, we place the origin at an arbitrary point $h$ inside the needle. The coordinate axes $x,y$ are to be horizontal, the former pointing north along the magnetic meridian and the latter east. We measure the coordinate $z$ upward. The action of terrestrial magnetism on the element $e$ thus contributes to $\mathrm{d}\Omega$ the term $Te\mathrm{d}x$.

Similarly, let the volume of the second, fixed needle be divided into infinitesimal elements. Let $X,Y,Z$ and the quantity of free magnetism $E$ correspond to an arbitrary element. Finally, let $r=\BareRoot{((X-x)^2+(Y-y)^2+(Z-z)^2)}$. Then the action of the element $E$ on the element $e$ contributes to the sum $\mathrm{d}\Omega$ the term $\frac{eE\mathrm{d}r}{r^n}$, if that action is assumed to be inversely proportional to the power $r^n$ of the distance $r$.

If we denote by $N$ the value of $u$ at which the thread is untwisted, the moment of the thread's torsional force can be expressed as $\theta(N-u)$. We can think of this force as acting tangentially at each end of the horizontal diameter of the thread at $G$, with the value $=\frac{\theta(N-u)}{D}$, where $D$ denotes that diameter. It is easy to see that this contributes to the sum $\mathrm{d}\Omega$ the term $\theta(N-u)\mathrm{d}u$.

The weight of the particles of the needle clearly contributes nothing to the sum $\mathrm{d}\Omega$, since $u$ is the only variable. We therefore have
\[
\mathrm{d}\Omega=\Sigma Te\mathrm{d}x+\Sigma\frac{eE\mathrm{d}r}{r^n}+\theta(N-u)\mathrm{d}u
\]
where the sum in the first term runs over all elements $e$ and the sum in the second over all pairs of individual $e$ and individual $E$. It is therefore clear that the condition for stable equilibrium is that
\[
\Omega=\Sigma Tex-\Sigma\frac{eE}{(n-1)r^{n-1}}-\tfrac12\theta(N-u)^2
\]
attain a maximum.

\SourceLeaf{106}
\Article{15}
For our purpose it is advisable always to arrange the experiments so that the magnetic axis of each needle is horizontal and the two needles are at nearly the same height. We will therefore restrict the subsequent calculations to these assumptions.

Let us refer the coordinates of the points of the first needle to axes fixed in that needle and still intersecting at the point $h$. Let the first axis point in the direction of the magnetic axis, the second lie horizontally to the right of the first and the third point vertically upward. Let the element $e$ have coordinates $a,b,c$ with respect to these axes. Likewise, let $A,B,C$ be the coordinates of the element $E$ with respect to similar axes fixed in the second needle and intersecting at a point $H$ in that needle. We assume that this point is chosen near the middle of the needle and at the same height as the point $h$.

If we were concerned with only one experiment, the position of the point $H$ would indeed be most conveniently determined by its distance from the point $h$ and the direction of the line joining them. For our purpose, however, we always need several experiments corresponding to different positions of the point $H$. These positions all lie on the same straight line, which need not pass exactly through the point $h$. It is therefore better to choose the notation from the outset so that a set of such experiments depends on a single variable. We will accordingly refer the point $H$ to an arbitrary point $h'$ in the same horizontal plane near $h$, with coordinates $\alpha,\beta,0$, and set the distance $h'H=R$. We set the angle between the line $h'H$ and the magnetic meridian $=\psi$. If we now denote the angle between the magnetic axis of the second needle and the magnetic meridian by $U$, we have
\[
\begin{aligned}
x&=a\cos u-b\sin u\\
y&=a\sin u+b\cos u\\
z&=c\\
X&=\alpha+R\cos\psi+A\cos U-B\sin U\\
Y&=\beta+R\sin\psi+A\sin U+B\cos U\\
Z&=C
\end{aligned}
\]

Everything is now ready for expanding the sum $\Omega$ and the quotient $\frac{\mathrm{d}\Omega}{\mathrm{d}u}$, which must vanish in equilibrium.

